%% !TEX root = manuscript.tex
\section*{Glossaire}
\vspace{1em}

\subsection*{Termes métier de l'immobilier}
\begin{longtable}{p{4cm} p{11cm}}
\textbf{Terme} & \textbf{Définition} \\
\hline
\endhead
\textbf{Acte de vente} & Document légal qui officialise la vente d'un bien immobilier entre un vendeur et un acquéreur. \\
\textbf{ADV} & Dossier de vente lié à un lot immobilier, ouvert lors d'une réservation et suivi jusqu'à l'acte de vente. Par extension, désigne dans nos extractions l'acquéreur principal du dossier. \\
\textbf{Aménagement foncier} & Opération qui consiste à diviser un terrain en plusieurs lots viabilisés, destinés à la vente pour la construction. \\
\textbf{Désistement} & Renoncement d'un acquéreur à une réservation avant l'acte de vente. \\
\textbf{Lot} & Terrain ou logement mis en vente, issu d'un découpage foncier ou d'un programme immobilier. \\
\textbf{Macro-lot} & Ensemble de lots vendus groupés, souvent destiné à un promoteur ou à un investisseur pour un projet plus large. \\
\textbf{Marge d'engagement} & Marge prévisionnelle d'une opération telle qu'elle est arrêtée au moment de l'engagement, et qui sert de référence pour mesurer la dérive ultérieure. \\
\textbf{Promotion immobilière} & Activité consistant à concevoir, financer, construire et commercialiser des logements ou bâtiments neufs. \\
\textbf{Réservation} & Engagement d'achat signé par un acquéreur en amont de l'acte de vente. C'est l'unité dans laquelle se mesure le rythme de commercialisation. \\
\textbf{Surface habitable} & Surface mesurée en mètres carrés, utilisée pour évaluer la taille d'un lot et estimer son prix. \\
\textbf{Tranche} & Subdivision d'une opération immobilière, distinguée selon qu'elle relève des travaux ou de la commercialisation. \\
\end{longtable}

\vspace{1em}
\subsection*{Outils et termes techniques}
\begin{longtable}{p{4cm} p{11cm}}
\textbf{Terme} & \textbf{Définition} \\
\hline
\endhead
\textbf{API} & \emph{Application Programming Interface}, interface de programmation. Elle expose les objets d'un logiciel de manière à pouvoir les créer, les lire ou les modifier par requêtes, sans passer par l'interface graphique. \\
\textbf{Coquilles vides} & Stratégie de déploiement consistant à créer d'abord la structure d'une opération, puis à en injecter le contenu budgétaire dans un second temps, afin que les utilisateurs puissent travailler sans attendre la reprise complète. \\
\textbf{CRM} & \emph{Customer Relationship Management}, outil de gestion de la relation client. Chez Angelotti, il s'agit de Vadimm. \\
\textbf{DirectAPI} & Module de l'éditeur Salvia permettant d'injecter dans l'ERP des fichiers tabulaires, sans passer par la saisie manuelle. \\
\textbf{EFR} & État Financier de Référence, l'un des exports budgétaires du système de gestion, qui donne pour chaque poste le budget, l'engagé et le facturé. \\
\textbf{ERP} & \emph{Enterprise Resource Planning}, progiciel de gestion intégré couvrant l'ensemble des processus de l'entreprise. \\
\textbf{GED} & Gestion Électronique des Documents. \\
\textbf{Grimmo} & ERP métier historique du groupe, remplacé cette année par SPO. \\
\textbf{JSON Patch} & Format normalisé décrivant la modification partielle d'un objet sous la forme d'opérations élémentaires, plutôt qu'en renvoyant l'objet entier. C'est lui qui a permis d'ajouter des coordonnées à des tiers déjà créés. \\
\textbf{MyReport} & Outil décisionnel du groupe, qui requête les bases de données sources et produit les tableaux de bord des services. \\
\textbf{Postman} & Client d'interface de programmation, permettant de composer des requêtes, de les organiser en collections et de les exécuter par lots. \\
\textbf{RGPD} & Règlement Général sur la Protection des Données, qui encadre le traitement des données à caractère personnel. \\
\textbf{Salvia} & Éditeur de logiciels spécialisés dans la gestion immobilière, dont SPO. \\
\textbf{SPO} & Suivi des Programmes Opérationnels, l'ERP métier qui remplace Grimmo depuis cette année. \\
\textbf{SQLite} & Moteur de base de données relationnelle léger, contenu dans un simple fichier, utilisé pour la base analytique du Copilote Financier. \\
\textbf{SSO} & \emph{Single Sign-On}, authentification unique donnant accès à plusieurs applications avec un seul compte. \\
\textbf{Streamlit} & Bibliothèque Python permettant de construire une application web interactive à partir d'un script. \\
\textbf{Tiers} & Dans SPO, entité désignant un client, une société ou un partenaire, à laquelle sont rattachés des rôles, des adresses et des coordonnées. \\
\textbf{Vadimm} & CRM du groupe, source des données commerciales reprises au second semestre. \\
\end{longtable}
